\documentclass[a4paper,12pt]{article}

\usepackage[french]{babel}
\usepackage[T1]{fontenc}
\usepackage[utf8]{inputenc}
\usepackage{amsmath, amssymb}
\usepackage{geometry}
\usepackage{graphicx}
\usepackage{array}
\usepackage{enumitem}
\usepackage{tkz-tab}
\usepackage{tikz}
\usepackage{pgfplots}
\pgfplotsset{compat=1.18}
\usepgfplotslibrary{statistics}

\geometry{top=2cm, bottom=2cm, left=2cm, right=2cm}

% Configuration
\newcommand{\sujet}{2}
\newcommand{\matiere}{Mathématiques}
\newcommand{\classe}{Première Générale spécialité maths}
\newcommand{\duree}{1 heure}
\newcommand{\datedevoir}{Mars 2026}
\newcommand{\theme}{Bac blanc -- Épreuve anticipée}
\newcommand{\totalpoints}{20}

% Cadre réponse pour les exercices
\newcommand{\reponse}{\par\vspace{0.5cm}
\framebox[\linewidth]{\begin{minipage}{\dimexpr\linewidth-2\fboxsep}
\vspace{2cm}\hspace{0.5cm}
\end{minipage}}\vspace{0.5cm}
}

\begin{document}

\begin{center}\Large\textbf{Épreuve anticipée de maths - Sujet \sujet}\end{center}

\begin{tabular}{|p{7.5cm}|p{7.5cm}|}
\hline
\textbf{Nom et Prénom :} & \textbf{Date : \datedevoir} \\
\framebox[7cm]{\hspace{0.1cm}\begin{minipage}{6.8cm}\vspace{1.5cm}\end{minipage}} & \\
\hline
\end{tabular}

\begin{center}\textbf{\classe\ -- \theme}\end{center}

\textbf{Durée : \duree\ -- Calculatrice interdite}

\section*{Partie 1 : Automatismes (12 questions -- 6 points)}
\textit{Pour chaque question, une seule réponse est correcte. 0,5 point par bonne réponse.}

\begin{enumerate}

\item On considère la relation \(M = x + \dfrac{y - z}{2t}\). Lorsque \(x = -1,\; y = 5,\; z = 1,\; t = -2\), la valeur de \(M\) est égale à :
\begin{itemize}
\item[A.] \(-2\)
\item[B.] \(0\)
\item[C.] \(2\)
\item[D.] \(-3\)
\end{itemize}

\item On considère trois fonctions définies sur \(\mathbb{R}\) :
\[k_1(x) = (2x-1)^2 - (x+1)(2x-1),\quad
k_2(x) = 3x - \left(5 - \frac{x}{2}\right),\quad
k_3(x) = \frac{4x^2 - (2x-3)^2}{2}\]
Parmi ces trois fonctions, lesquelles sont des fonctions affines ?
\begin{itemize}
\item[A.] aucune
\item[B.] toutes
\item[C.] uniquement \(k_2\)
\item[D.] uniquement \(k_2\) et \(k_3\)
\end{itemize}

\item Voici une série de notes avec coefficients :
\[
\begin{array}{|c|c|c|c|c|}
\hline
\text{Note} & 12 & 8 & 15 & a \\
\hline
\text{Coefficient} & 1 & 1 & 2 & 1 \\
\hline
\end{array}
\]
On note \(m\) la moyenne. Que doit valoir \(a\) pour que \(m = 12\) ?
\begin{itemize}
\item[A.] \(7\)
\item[B.] \(1\)
\item[C.] \(10\)
\item[D.] \(4{,}25\)
\end{itemize}

\item Un boulanger observe que sur 260 viennoiseries vendues, 85 sont des pains au chocolat. Les pains au chocolat représentent alors (en pourcentage des ventes, arrondi à l'unité) :
\begin{itemize}
\item[A.] \(25\%\)
\item[B.] \(33\%\)
\item[C.] \(75\%\)
\item[D.] \(35\%\)
\end{itemize}

\item On a représenté ci-dessous les boîtes à moustaches de deux séries statistiques.
\begin{center}
\begin{tikzpicture}
\begin{axis}[
  title={Série 1},
  xlabel={Valeurs},
  ytick={1}, yticklabels={},
  axis y line=none,
  axis x line=bottom,
  xmin=0, xmax=20,
  xtick={0,1,...,20},
  width=12cm, height=3cm,
]
\addplot+[boxplot prepared={
  lower whisker=5, lower quartile=7,
  median=9, upper quartile=13,
  upper whisker=15
}] coordinates {};
\end{axis}
\end{tikzpicture}

\vspace{0.5cm}

\begin{tikzpicture}
\begin{axis}[
  title={Série 2},
  xlabel={Valeurs},
  ytick={1}, yticklabels={},
  axis y line=none,
  axis x line=bottom,
  xmin=0, xmax=20,
  xtick={0,1,...,20},
  width=12cm, height=3cm,
]
\addplot+[boxplot prepared={
  lower whisker=3, lower quartile=7,
  median=9, upper quartile=13,
  upper whisker=19
}] coordinates {};
\end{axis}
\end{tikzpicture}
\end{center}
On peut affirmer que :
\begin{itemize}
\item[A.] Dans la série 1, la médiane est égale à 9.
\item[B.] Dans la série 2, au moins \(75\%\) des valeurs sont inférieures ou égales à 13.
\item[C.] Les deux séries ont la même étendue.
\item[D.] L'écart interquartile est plus grand pour la série 1 que pour la série 2.
\end{itemize}

\item Aide au calcul : \(\pi \approx 3{,}14\). L'ordre de grandeur de \((3\pi)^2 \times 120\) est :
\begin{itemize}
\item[A.] \(1000\)
\item[B.] \(10\,000\)
\item[C.] \(100\,000\)
\item[D.] \(1\,000\,000\)
\end{itemize}

\item \(\displaystyle \frac{4^3 \times 3^2}{6^2} =\)
\begin{itemize}
\item[A.] \(4\)
\item[B.] \(8\)
\item[C.] \(16\)
\item[D.] \(2\)
\end{itemize}

\item La fonction \(f\) définie sur \(\mathbb{R}\) par \(f(x) = 3x - 1\) admet pour tableau de signe :
\begin{itemize}
\item[A.]
$\begin{array}{c|ccc}
x & -\infty & \frac{1}{3} & +\infty \\ \hline
f(x) & - & 0 & +
\end{array}$
\item[B.]
$\begin{array}{c|ccc}
x & -\infty & \frac{1}{3} & +\infty \\ \hline
f(x) & + & 0 & -
\end{array}$
\item[C.]
$\begin{array}{c|ccc}
x & -\infty & 0 & +\infty \\ \hline
f(x) & - & 0 & +
\end{array}$
\item[D.]
$\begin{array}{c|ccc}
x & -\infty & 0 & +\infty \\ \hline
f(x) & + & 0 & -
\end{array}$
\end{itemize}

\item Une bactérie de \(0{,}8\)\,g voit sa masse diminuer de \(25\%\). La nouvelle masse est égale à :
\begin{itemize}
\item[A.] \(0{,}6\)\,g
\item[B.] \(1\)\,g
\item[C.] \(0{,}75\)\,g
\item[D.] \(0{,}2\)\,g
\end{itemize}

\item La loi d'interaction gravitationnelle entre deux corps \(C_1\) et \(C_2\) est
\(F = G\dfrac{m_1 m_2}{d^2}\) où \(m_1, m_2\) sont les masses, \(d\) la distance et \(G\) la constante gravitationnelle. On a :
\begin{itemize}
\item[A.] \(d = \dfrac{F - G}{2m_1 m_2}\)
\item[B.] \(d = \sqrt{\dfrac{G m_1 m_2}{F}}\)
\item[C.] \(d = \dfrac{G m_1 m_2}{F^2}\)
\item[D.] \(d = \sqrt{\dfrac{F}{G m_1 m_2}}\)
\end{itemize}

\item L'expression développée réduite de \((5x+4)^2\) est :
\begin{itemize}
\item[A.] \(25x^2 + 8x + 16\)
\item[B.] \(10x^2 + 20x + 16\)
\item[C.] \(25x^2 + 16x + 4\)
\item[D.] \(25x^2 + 40x + 16\)
\end{itemize}

\item On considère un groupe d'élèves dont la répartition suivant la LV2 choisie et le sexe est donnée ci-dessous :
\[
\begin{array}{|c|c|c|c|}
\hline
& \text{Allemand} & \text{Espagnol} & \text{Total} \\
\hline
\text{Filles} & 30 & 20 & 50 \\
\hline
\text{Garçons} & 40 & 30 & 70 \\
\hline
\text{Total} & 70 & 50 & 120 \\
\hline
\end{array}
\]
On choisit un élève au hasard. Quelle est la probabilité que l'élève suive allemand sachant que c'est une fille ?
\begin{itemize}
\item[A.] \(0{,}25\)
\item[B.] \(0{,}6\)
\item[C.] \(\dfrac{3}{7}\)
\item[D.] \(\dfrac{4}{7}\)
\end{itemize}

\end{enumerate}

\section*{Partie 2 : Exercices (14 points)}

\subsection*{Exercice 1 (7 points)}
À partir du printemps 2026, Lina gère un bassin d'irrigation contenant au départ \(1000\ \mathrm{m}^3\) d'eau. Chaque année, à cause de l'évaporation et des fuites, \(20\%\) du volume d'eau est perdu. En fin d'été, Lina rajoute \(40\ \mathrm{m}^3\) d'eau dans le bassin.

Pour tout entier naturel \(n\), on note \(u_n\) le volume d'eau (en \(\mathrm{m}^3\)) dans le bassin au début du printemps de l'année \(2026+n\). Ainsi \(u_0 = 1000\).

\begin{enumerate}
\item Calculer le volume d'eau du bassin au printemps 2027 suivant ce modèle. \reponse
\item Justifier que, pour tout entier naturel \(n\), \(u_{n+1} = 0{,}8\,u_n + 40\). \reponse
\item Justifier que la suite \((u_n)\) n'est pas arithmétique. \reponse
\item On considère la suite \((v_n)\) définie pour tout entier naturel \(n\) par \(v_n = u_n - 200\).
\begin{enumerate}
\item Montrer que la suite \((v_n)\) est géométrique de raison \(0{,}8\). \reponse
\item En déduire l'expression de \(v_n\) en fonction de \(n\). \reponse
\item En déduire que, pour tout entier naturel \(n\), \(u_n = 800 \times 0{,}8^n + 200\). \reponse
\end{enumerate}
\item À l'aide d'un tableur, on a obtenu les valeurs arrondies au centième suivantes :
\[
\begin{array}{|c|c|}
\hline
n & u_n \\
\hline
5 & 409{,}83 \\
10 & 273{,}09 \\
15 & 228{,}36 \\
20 & 206{,}59 \\
25 & 201{,}05 \\
30 & 200{,}17 \\
\hline
\end{array}
\]
\begin{enumerate}
\item Conjecturer le comportement des termes de \((u_n)\) lorsque \(n\) tend vers \(+\infty\). \reponse
\item Interpréter ce comportement dans le contexte de l'exercice. \reponse
\end{enumerate}
\end{enumerate}

\subsection*{Exercice 2 (7 points)}
Soit \(f\) la fonction définie pour tout nombre réel \(x \neq -\frac{1}{2}\) par
\[
f(x) = \frac{3x - 3}{2x + 1}.
\]

\begin{enumerate}
\item On admet que \(f\) est dérivable sur \(\mathbb{R}\setminus\left\{-\frac{1}{2}\right\}\). Montrer que pour tout \(x \neq -\frac{1}{2}\),
\[
f'(x) = \frac{9}{(2x+1)^2}.
\] \reponse
\item En déduire les variations de \(f\) sur son ensemble de définition. \reponse
\item On note \(\mathcal{C}\) la courbe représentative de \(f\). 
\begin{enumerate}
\item Montrer que \(\mathcal{C}\) admet exactement deux tangentes de coefficients directeurs égaux à \(1\).
\textit{Indication : Résoudre \(f'(x) = 1\).} \reponse
\item Déterminer l'équation réduite de chacune de ces tangentes. \reponse
\end{enumerate}
\end{enumerate}

\vspace{1cm}
\textbf{Barème indicatif :}
\begin{itemize}
\item Automatismes : 6 points (0,5 pt par question)
\item Exercice 1 : 7 points
\item Exercice 2 : 7 points
\end{itemize}

\end{document}